% =====================================================================
%  APÊNDICE C — CONFIGURAÇÃO DO AMBIENTE
% =====================================================================

\chapter{CONFIGURAÇÃO DO AMBIENTE}
\label{ap:ambiente}

Este apêndice descreve a configuração dos principais componentes do ambiente de
execução do \argus{}.

\section{Conexão com sensores Bluetooth Low Energy}
\label{ap:ambiente-ble}

Os sensores BLE (sensor de frequência cardíaca, medidor de potência e rolo
inteligente) são detectados automaticamente pelo escaneamento do sistema. No
Windows, caso a
pilha Bluetooth não exponha os UUIDs de serviço durante o escaneamento, o
\argus{} lista os dispositivos nomeados e valida os serviços FTMS/HR após a
seleção e conexão. Para diagnóstico, a variável de ambiente
\code{ARGUS\_BLE\_DEBUG=1} ativa logs detalhados de escaneamento.

O pareamento é feito pela interface de conexão do \argus{}. No Windows, o
sistema operacional precisa permitir o acesso ao Bluetooth nas configurações de
privacidade; no Linux, o usuário deve ter permissão para acessar o adaptador
Bluetooth; no Android, a permissão de localização/dispositivos próximos deve ser
concedida ao aplicativo. O rolo deve estar livre de outro aplicativo no
momento da conexão, porque alguns modelos aceitam apenas um cliente BLE por vez.

\section{Assistente de treino com IA local (Ollama)}
\label{ap:ambiente-ollama}

O assistente de IA requer uma instância do Ollama \cite{ollama2024} executando
no dispositivo. Recomenda-se o modelo quantizado de pequena escala
\texttt{qwen2.5:3b}, que executa de forma responsiva em CPU. A seleção do
modelo é feita na interface do \argus{} (\textit{AIChatController}).

A instalação é feita a partir de \url{https://ollama.com/}. Depois de instalar
o serviço, baixe o modelo recomendado:

\begin{verbatim}
ollama pull qwen2.5:3b
\end{verbatim}

Por padrão, o \argus{} procura o serviço em \code{http://localhost:11434}. Para
usar outro \textit{endpoint}, configure a variável \code{OLLAMA\_HOST} no
ambiente do serviço e ajuste a URL na interface de configuração do assistente.
Modelos maiores podem ser usados, mas exigem mais memória e tornam a inferência
em CPU mais lenta.

\section{Integração com o Strava e o Intervals.icu}
\label{ap:ambiente-strava}

As integrações com o Strava e o Intervals.icu exigem o cadastro de uma
aplicação no respectivo portal de desenvolvedores e o armazenamento das
credenciais de OAuth 2.0. O fluxo de autorização é realizado por um servidor
HTTP local, e o \textit{refresh token} é gerenciado automaticamente pelo
sistema.

Para o Strava, é necessário criar uma aplicação em
\url{https://developers.strava.com/} e obter \code{client\_id} e
\code{client\_secret}. A URL de redirecionamento deve apontar para o servidor
local usado pelo \argus{} durante o fluxo OAuth. As permissões solicitadas
incluem leitura do perfil e escrita de atividades. Para o Intervals.icu, o
usuário gera uma chave de API nas configurações da conta e a informa na
interface de integrações. As credenciais não devem ser versionadas no
repositório; em produção, devem ser armazenadas de forma protegida no
dispositivo.
